% ============================================================================
% Fichier  : partie2/P02_C07_decomposer_coter_conclure.tex
% Titre    : Décomposer, Coter, Conclure
% Auteur   : MINKA MI NGUIDJOÏ Thierry Emmanuel
% Version  : 1.0 -- novembre 2026
% Corrigés : corriges/Corr_C07.tex
% ============================================================================
% Méthode reprise sans modification du corpus (ACSI, chapitres 7 et 9) :
% gabarit canonique à huit niveaux (article A0) ; fiche d'assertion à cinq
% attributs ; règle d'atomicité ; trois conditions (exhaustivité,
% non-redondance, borne cognitive) ; cote d'un contrôle à plusieurs critères
% (un NON suffit ; à défaut un NEK ; OUI si tous OUI) ; mesures mu, U, H
% (article A1) ; quadrant aux seuils du cours (mu 0,5 ; U 0,20) ; règle de
% décision « U d'abord, mu ensuite, H pour contrôler, les cotes de
% l'auditeur prévalent » ; conversion en maturité (convention ACSI) ;
% observation en sept éléments. Domaine D de LiZbiZ : valeurs d'ACSI.
% Faits du cas (bible) : mission de mars 2028 sur le réseau d'agences, T1 à
% T8 = 3/3/2. Faits nouveaux : décomposition en six sous-domaines S1 à S6 ;
% exercice inopiné de Kribi le 14 mars 2028 (2 h 15), contradictoire du
% 20 mars ; risque R-AG-02 (compromission d'un GAB).
% ISO/IEC 27002:2022, mesure 5.30 (préparation des TIC pour la continuité
% d'activité) : citée par son numéro et son intitulé.
% ============================================================================

\chapter{Décomposer, Coter, Conclure}
\label{chap:conclure}
\chaptermark{Décomposer et conclure}

\epigraph{%
    « Conclure, ce n'est pas compter les \emph{oui}.\\
    C'est savoir ce qu'ils valent. »%
}{Principe de la méthode}

\niveauChapitre{Intermédiaire}{%
    Maîtrise du Chapitre~\ref{chap:audit}%
}

\begin{boxobjectifs}
\begin{itemize}
    \item Décomposer un domaine de sécurité selon le gabarit canonique à
          huit niveaux, et vérifier la qualité de la décomposition.
    \item Rédiger une fiche d'assertion à plusieurs critères, en respectant
          l'atomicité.
    \item Déduire la cote d'un contrôle de celles de ses critères.
    \item Calculer $\mu$, $U$ et $H$ par sous-domaine, et localiser
          l'incertitude.
    \item Situer un domaine dans le quadrant conformité~$\times$ incertitude
          et appliquer la règle de décision.
    \item Rédiger une observation en sept éléments et la conclusion d'un
          rapport.
\end{itemize}
\end{boxobjectifs}

\bigskip

% ============================================================================
\section*{Avant de commencer : « Alors, c'est sécurisé~? »}
% ============================================================================

Douala, avril 2028. Mme~EKAMBI présente les résultats de la mission sur le
réseau d'agences au comité des risques~: huit tests, 3~\OUI{}, 3~\NON{},
2~\NEK{} (Chapitre~\ref{chap:audit}). Un administrateur l'interrompt~:
« Trois sur huit, c'est 37~\%. Les agences sont-elles sécurisées, oui ou
non~? »

La question est légitime, la réponse ne peut pas être un pourcentage. Les
huit cotes ne portent pas sur la même chose~: trois concernent les
barrières techniques du scénario R-AG-01, une la capacité à continuer de
servir, une le personnel, une un prestataire qui a refusé l'accès. Les
additionner, c'est mélanger un pare-feu et une limitation de l'étendue.

\medskip

Ce chapitre apprend à structurer ce que l'on audite, à coter chaque
contrôle selon ses critères, puis à conclure, en disant à la fois ce qui va
mal et ce que l'on ne sait pas. La méthode est celle de l'auditologie,
commune aux cours du corpus~\cite{minka2026auditologie,minka2026nek}.

\begin{boxalert}
\textbf{Violation typique.} Résumer un audit par le pourcentage de
contrôles conformes.

\smallskip
\textbf{Règle~:} une conclusion se donne par sous-domaine, avec trois
mesures (conformité, incertitude, dispersion), selon une règle de décision
écrite à l'avance. Un taux unique cache où sont les défauts et où est
l'ignorance.
\end{boxalert}

\separateur

% ============================================================================
\section{Décomposer}
\label{sec:c07-decomposer}
% ============================================================================

\subsection{Pourquoi décomposer}

Un domaine de sécurité compte des dizaines de contrôles. La mémoire de
travail humaine traite environ sept éléments à la fois~\cite{miller1956}~:
au-delà, l'auditeur oublie, recouvre ou mélange. La décomposition découpe
l'objet audité en niveaux emboîtés, chacun de taille maîtrisable. Elle a
trois usages~: vérifier que rien n'échappe au programme, situer chaque
cote, et localiser l'incertitude.

\subsection{Le gabarit canonique}

Le programme de l'auditologie propose un \termecle{gabarit
canonique}\index{gabarit canonique} à huit niveaux~\cite{minka2026auditologie}.

\begin{tableaucours}{cL{2.2cm}L{4.4cm}L{4.6cm}}{%
    \tableauHeader{N°} &
    \tableauHeader{Niveau} &
    \tableauHeader{Question} &
    \tableauHeader{Exemple (réseau d'agences)}
}
1 & Domaine & Que soumet-on à l'audit~? & Sécurité du réseau d'agences \\
2 & Sous-domaine & Quels grands ensembles le composent~? & Récupération \\
3 & Processus & Quelles suites d'activités~? & Continuité du service au
guichet \\
4 & Activité & Quelles actions élémentaires~? & Servir en mode dégradé \\
5 & Contrôle & Quel dispositif maîtrise l'activité~? & Procédure de mode
dégradé, exercée \\
6 & Test & Comment l'auditeur le vérifie-t-il~? & Exercice inopiné,
chronométré \\
7 & Preuve & Quel élément le test produit-il~? & Compte rendu signé,
chronométrage \\
8 & Trace & Quel enregistrement l'atteste~? & Journal des opérations en
mode dégradé \\
\end{tableaucours}
\vspace{-0.8em}
\begin{center}{\small\itshape Tableau~7.1. Le gabarit canonique
(d'après~\cite{minka2026auditologie})}\end{center}

Les quatre premiers niveaux disent \emph{ce qui est audité}~; les quatre
derniers, \emph{comment on le vérifie}. Le contrôle fait la charnière~:
c'est l'objet sur lequel l'auditeur se prononce.

\subsection{Trois conditions de qualité}

\begin{tableaucours}{L{3.0cm}L{5.2cm}L{4.0cm}}{%
    \tableauHeader{Condition} &
    \tableauHeader{Sens} &
    \tableauHeader{Défaut si elle manque}
}
Exhaustivité & Les éléments d'un niveau couvrent entièrement leur parent &
Une partie du domaine n'est pas auditée \\
Non-redondance & Deux éléments d'un même niveau ne se recouvrent pas~;
chacun a un seul parent & Double comptage, conclusions contradictoires \\
Borne cognitive & Aucun élément n'a plus de sept enfants environ & Grille
illisible, éléments oubliés \\
\end{tableaucours}
\vspace{-0.8em}
\begin{center}{\small\itshape Tableau~7.2. Les trois conditions d'une bonne
décomposition}\end{center}

\subsection{Le réseau d'agences, décomposé}

\begin{boxscenario}{Opération SANAGA. Six sous-domaines}
Mme~EKAMBI décompose le domaine « sécurité du réseau d'agences » en six
sous-domaines. Les quatre premiers suivent les étapes du scénario R-AG-01
(Chapitre~\ref{chap:risque})~; les deux derniers couvrent ce que le
scénario ne prévoyait pas.

\medskip
\begin{tableaucours}{cL{3.4cm}L{4.8cm}c}{%
    \tableauHeader{Code} &
    \tableauHeader{Sous-domaine} &
    \tableauHeader{Contrôles testés} &
    \tableauHeader{Tests}
}
S1 & Entrée (accès, authentification) & Second facteur sur tous les
comptes & T1 \\
S2 & Propagation (cloisonnement, privilèges) & Cloisonnement par agence~;
fin de l'administration locale & T2, T7 \\
S3 & Détection & Journalisation centralisée & T6 \\
S4 & Récupération & Sauvegardes hors ligne testées~; mode dégradé & T3,
T4 \\
S5 & Facteur humain & Sensibilisation des agents & T5 \\
S6 & Équipements de tiers & Correctifs des GAB & T8 \\
\end{tableaucours}

\medskip
\textbf{Vérification.} Six sous-domaines, sous la borne de sept~;
chaque test a un seul parent~; l'ensemble couvre le périmètre de la lettre
de mission. L'exhaustivité appelle une réserve~: d'autres contrôles
existent dans chaque sous-domaine (par exemple les correctifs des postes en
S1)~; la mission n'a testé que ceux du plan de traitement. Le rapport doit
le dire.
\end{boxscenario}

% ============================================================================
\section{Coter un Contrôle}
\label{sec:c07-coter}
% ============================================================================

\subsection{La fiche d'assertion}

Le gabarit dit \emph{où} se trouve un contrôle. La \termecle{fiche
d'assertion}\index{fiche d'assertion} dit ce que l'auditeur affirme à son
sujet, et ce qui arrive si l'affirmation est fausse.

\begin{tableaucours}{L{2.4cm}L{5.0cm}L{4.8cm}}{%
    \tableauHeader{Attribut} &
    \tableauHeader{Rôle} &
    \tableauHeader{Forme}
}
Objectif & Ce que le contrôle doit garantir & À l'infinitif~: « S'assurer
que\ldots » \\
Critère & La condition vérifiable qui définit la conformité & Au présent~:
une condition unique, observable \\
Question & La question posée pour coter & Interrogative~; réponse \OUI{},
\NON{} ou \NEK{} \\
Risque & L'état de non-conformité si la réponse est \NON{} & L'antithèse du
critère \\
Conséquence & L'effet sur l'organisation si le risque se réalise &
Financier, opérationnel, réglementaire \\
\end{tableaucours}
\vspace{-0.8em}
\begin{center}{\small\itshape Tableau~7.3. Les attributs de la fiche
d'assertion}\end{center}

\begin{boxinfo}
\textbf{Règle d'atomicité.} L'objectif, le critère et la question ne
contiennent aucun «~et~» ni «~ou~». Toute assertion qui en contient est
scindée en autant d'assertions qu'il y a de conditions indépendantes. Ainsi
une réponse \NON{} désigne un seul défaut.
\end{boxinfo}

\subsection{Un contrôle, trois critères}

\begin{boxscenario}{Opération SANAGA. Fiche d'assertion du contrôle S4.2}
\textbf{Position~:} domaine «~réseau d'agences~» $\rightarrow$ S4
(récupération) $\rightarrow$ continuité du service au guichet
$\rightarrow$ servir en mode dégradé $\rightarrow$ contrôle S4.2.

\medskip
\textbf{Objectif~:} s'assurer que chaque agence peut servir ses clients
sans le système d'information habituel.

\medskip
\textbf{Critère~1~:} la procédure de mode dégradé est disponible sur place,
sous forme papier. \emph{Question~:} la procédure est-elle présente dans
l'agence~?

\textbf{Critère~2~:} les codes d'accès au mode dégradé sont connus de la
personne qui dirige l'agence le jour du test. \emph{Question~:} le
responsable présent connaît-il les codes~?

\textbf{Critère~3~:} l'agence sert un premier client en moins d'une heure.
\emph{Question~:} le premier client a-t-il été servi en moins d'une
heure~?

\medskip
\textbf{Test, preuve, trace~:} exercice inopiné dans deux agences de
l'échantillon~; compte rendu signé et chronométrage~; journal des
opérations en mode dégradé.

\medskip
\textbf{Risque~:} une agence touchée par un incident reste fermée.

\textbf{Conséquence~:} perte d'exploitation (190~millions pour neuf agences
en 2026, Chapitre~\ref{chap:fair}), atteinte à la réputation.
\end{boxscenario}

\subsection{De la cote des critères à celle du contrôle}

Chaque critère reçoit sa propre cote. La cote du contrôle s'en déduit par
une règle simple~:
\begin{itemize}
    \item un seul critère \NON{} suffit pour coter le contrôle \NON{}~;
    \item à défaut, un seul critère \NEK{} donne \NEK{}~;
    \item le contrôle n'est coté \OUI{} que si tous ses critères le sont.
\end{itemize}

À Garoua, les trois critères sont \OUI{}. À Kribi, le critère~1 est \NON{}
(procédure introuvable), le critère~2 \NON{} (codes inconnus du chef
d'agence remplaçant), donc le critère~3 aussi (2~h~15). Un contrôle
d'agence se teste agence par agence~: S4.2 est \NON{} pour l'échantillon,
puisqu'une agence sur deux ne le satisfait pas.

\begin{boxalert}
\textbf{Violation typique.} Coter « \OUI{} à 50~\% » un contrôle conforme
dans une agence et défaillant dans l'autre.

\smallskip
\textbf{Règle~:} la cote porte sur le critère, qui est binaire. On cote
\NON{}, et le papier de travail précise où le contrôle fonctionne et où il
échoue.
\end{boxalert}

% ============================================================================
\section{Mesurer et Localiser l'Incertitude}
\label{sec:c07-mesurer}
% ============================================================================

\subsection{Trois mesures}

\begin{boxdef}
Pour un ensemble de cotes, on note $p(\OUI)$, $p(\NON)$ et $p(\NEK)$ les
parts de chaque réponse~\cite{minka2026nek}.

\medskip
\termecle{Taux de conformité}\index{taux de conformité}~: $\mu = p(\OUI)$.

\medskip
\termecle{Taux d'incertitude}\index{taux d'incertitude}~: $U = p(\NEK)$.

\medskip
\termecle{Entropie}\index{entropie}~: dispersion des réponses, en bits,
\[
    H = -\sum_{m \in \{\OUI, \NON, \NEK\}} p(m) \log_2 p(m),
    \qquad 0 \leq H \leq \log_2 3 \approx 1{,}585,
\]
avec la convention $0 \log_2 0 = 0$.
\end{boxdef}

\begin{tableaucours}{L{4.6cm}cccL{3.4cm}}{%
    \tableauHeader{Situation} &
    \tableauHeader{$\mu$} & \tableauHeader{$U$} & \tableauHeader{$H$} &
    \tableauHeader{Lecture}
}
Tout est \OUI{} & 1 & 0 & 0 & Conformité certaine \\
Tout est \NON{} & 0 & 0 & 0 & Non-conformité certaine \\
Tout est \NEK{} & 0 & 1 & 0 & \textbf{Ignorance unanime} \\
Réponses également réparties & $1/3$ & $1/3$ & 1,585 & Dispersion maximale \\
\end{tableaucours}
\vspace{-0.8em}
\begin{center}{\small\itshape Tableau~7.4. Les cas limites des trois
mesures (d'après~\cite{minka2026nek})}\end{center}

\begin{boxalert}
\textbf{Violation typique.} Conclure « certitude » parce que l'entropie
est nulle.

\smallskip
\textbf{Règle~:} une entropie nulle dit seulement que toutes les cotes sont
identiques. Si elles sont toutes \NEK{}, on ne sait rien~: $H = 0$ et
$U = 1$. On ne lit jamais $H$ sans $U$ ni $\mu$.
\end{boxalert}

\subsection{Par sous-domaine}

\begin{boxlecture}{Le réseau d'agences, sous-domaine par sous-domaine}
\begin{tableaucours}{cL{3.4cm}cccc}{%
    \tableauHeader{Code} &
    \tableauHeader{Sous-domaine} &
    \tableauHeader{\OUI{}/\NON{}/\NEK{}} &
    \tableauHeader{$\mu$} & \tableauHeader{$U$} & \tableauHeader{$H$}
}
S1 & Entrée & 1 / 0 / 0 & 1,00 & 0,00 & 0,00 \\
S2 & Propagation & 1 / 1 / 0 & 0,50 & 0,00 & 1,00 \\
S3 & Détection & 0 / 1 / 0 & 0,00 & 0,00 & 0,00 \\
S4 & Récupération & 1 / 1 / 0 & 0,50 & 0,00 & 1,00 \\
S5 & Facteur humain & 0 / 0 / 1 & 0,00 & 1,00 & 0,00 \\
S6 & Équipements de tiers & 0 / 0 / 1 & 0,00 & 1,00 & 0,00 \\
\textbf{Domaine} & & \textbf{3 / 3 / 2} & \textbf{0,375} & \textbf{0,25} &
\textbf{1,56} \\
\end{tableaucours}

\medskip
\textbf{Constatation.} Toute l'incertitude du domaine se concentre dans S5
et S6~: deux cas d'ignorance unanime ($H = 0$, $U = 1$). Les quatre autres
sous-domaines sont connus ($U = 0$)~: S1 est conforme, S3 non conforme, S2
et S4 partagés.

\medskip
\textbf{Analyse} (confiance élevée). Lu d'un bloc, le domaine paraît
« à 37~\% » et en partie inconnu. Décomposé, il dit trois choses
distinctes~: les barrières techniques du scénario R-AG-01 tiennent (S1,
le cloisonnement de S2, les sauvegardes de S4)~; des défauts connus restent
à corriger (administration locale, journalisation, mode dégradé)~; et la
banque ignore tout de deux sujets, le personnel et les GAB. Ce ne sont pas
les mêmes décisions.

\medskip
\hfill\lecture{0{,}375}{0{,}25}{1{,}56}
\end{boxlecture}

\begin{boxinfo}
\textbf{Peu de cotes, mesures fragiles.} S2 et S4 ne comptent que deux
contrôles chacun~: $\mu = 0{,}50$ y signifie « un sur deux », pas une
proportion stable. À ce niveau, on conclut contrôle par contrôle~; les
mesures servent surtout à comparer les sous-domaines et à repérer les
\NEK{}.
\end{boxinfo}

% ============================================================================
\section{Conclure}
\label{sec:c07-conclure}
% ============================================================================

\subsection{Le quadrant conformité~$\times$ incertitude}

L'article sur l'incertitude épistémique propose de lire ensemble la
conformité et l'incertitude, dans un quadrant à quatre
régions~\cite{minka2026nek}. Le cours retient, par convention, un seuil de
0,5 pour $\mu$ et de 0,20 pour $U$~; chaque mission peut fixer les siens, à
condition de les écrire dans le plan d'audit.

\begin{figure}[H]
\centering
\begin{tikzpicture}[x=6cm, y=4.5cm]
    \fill[red!15] (0,0) rectangle (0.5,0.2);
    \fill[green!18] (0.5,0) rectangle (1,0.2);
    \fill[orange!15] (0,0.2) rectangle (0.5,1);
    \fill[orange!25] (0.5,0.2) rectangle (1,1);
    \draw[thick, color=CouleurPrimaire] (0,0) rectangle (1,1);
    \draw[dashed, color=CouleurPrimaire] (0.5,0) -- (0.5,1);
    \draw[dashed, color=CouleurPrimaire] (0,0.2) -- (1,0.2);
    \node[align=center, font=\scriptsize\bfseries] at (0.25,0.1)
        {Non-conformité claire};
    \node[align=center, font=\scriptsize\bfseries] at (0.75,0.1)
        {Conformité claire};
    \node[align=center, font=\scriptsize\bfseries] at (0.25,0.6)
        {Non-conformité\\opaque};
    \node[align=center, font=\scriptsize\bfseries] at (0.75,0.6)
        {Conformité\\opaque};
    \fill[CouleurAccent] (0.375,0.25) circle (0.012);
    \node[font=\scriptsize, anchor=west] at (0.38,0.27) {Domaine};
    \fill[CouleurPrimaire] (0.02,0.98) circle (0.010);
    \node[font=\scriptsize, anchor=west] at (0.03,0.94) {S5, S6};
    \fill[CouleurPrimaire] (0.98,0.02) circle (0.010);
    \node[font=\scriptsize, anchor=east] at (0.97,0.045) {S1};
    \node[font=\scriptsize] at (0.5,-0.06) {0,5};
    \node[font=\scriptsize] at (-0.06,0.2) {0,20};
    \node[font=\small] at (0.5,-0.16) {Taux de conformité $\mu$};
    \node[font=\small, rotate=90] at (-0.15,0.5) {Taux d'incertitude $U$};
\end{tikzpicture}
\caption{Le quadrant diagnostique (d'après l'article sur l'incertitude
épistémique~; seuils du cours)}
\label{fig:c07-quadrant}
\end{figure}

\begin{tableaucours}{L{3.4cm}L{4.0cm}L{4.8cm}}{%
    \tableauHeader{Région} &
    \tableauHeader{Situation} &
    \tableauHeader{Action}
}
Conformité claire & Conformité établie, peu d'inconnu & Validation
courante \\
Non-conformité claire & Non-conformité établie, peu d'inconnu &
Correctifs \\
\textbf{Conformité opaque} & Beaucoup de \OUI{}, mais une part importante
de \NEK{} & \textbf{Alerte prioritaire}~: rechercher la preuve avant toute
conclusion \\
Non-conformité opaque & Beaucoup de \NON{} et une part importante de
\NEK{} & Rechercher la preuve et engager des premiers correctifs \\
\end{tableaucours}
\vspace{-0.8em}
\begin{center}{\small\itshape Tableau~7.5. Les quatre régions et leurs
actions (d'après~\cite{minka2026nek})}\end{center}

La conformité opaque est la région la plus dangereuse~: une grille binaire
la confond avec la conformité claire. Le rapport de la RSSI de février
2028 en donne l'exemple~: huit mesures déclarées « en place », soit une
conformité apparente de 100~\%~; mais ces \OUI{} reposaient sur des
déclarations, qui ne valaient que des \NEK{} tant qu'on ne les avait pas
vérifiées.

\subsection{La règle de décision}

\begin{boxprocedure}
\textbf{Conclure sur un domaine.} L'auditeur ne conclut qu'après avoir
consigné, pour chaque domaine et chaque sous-domaine~:
\begin{enumerate}
    \item \textbf{$U$ d'abord.} Si $U \geq 0{,}20$, l'ensemble est opaque~:
          aucune conclusion de conformité n'est possible. L'auditeur
          recherche la preuve, ou conclut \emph{sur l'opacité elle-même}.
    \item \textbf{$\mu$ ensuite.} Si $U < 0{,}20$, l'ensemble est clair~:
          conforme si $\mu \geq 0{,}5$, non conforme sinon.
    \item \textbf{$H$ pour contrôler.} Une entropie élevée avec $U$ faible
          signale un désaccord entre \OUI{} et \NON{}~: l'ensemble n'est pas
          homogène. On le décompose avant de conclure.
    \item \textbf{Les cotes de l'auditeur prévalent.} Les mesures orientent~;
          la conclusion repose sur les contrôles cotés, preuve à l'appui.
\end{enumerate}
\end{boxprocedure}

\begin{boxscenario}{Opération SANAGA. La conclusion de Mme EKAMBI}
\textbf{Domaine.} $U = 0{,}25 \geq 0{,}20$~: le domaine est opaque, en
région de non-conformité opaque ($\mu = 0{,}375$). Aucune conclusion de
conformité d'ensemble n'est possible~; la mission recherche les preuves
manquantes et engage des premiers correctifs.

\medskip
\textbf{Sous-domaines.}
\begin{itemize}
    \item S1~: clair et conforme.
    \item S2 et S4~: clairs, $H = 1$~: non homogènes~; conclusion par
          contrôle~: cloisonnement et sauvegardes efficaces, administration
          locale et mode dégradé non efficaces.
    \item S3~: clair, non conforme (corrigé en cours de mission).
    \item S5 et S6~: ignorance unanime~; l'auditeur conclut sur l'opacité.
          Pour S6, l'opacité vient d'une limitation de l'étendue, portée au
          comité.
\end{itemize}

\textbf{Phrase pour le comité.} « Les barrières techniques qui ont été
défaillantes en 2026 sont désormais en place et fonctionnent. Trois
défauts connus restent à corriger. Et la banque ne sait pas, à ce jour, si
ses agents sont formés ni si ses GAB sont protégés~: sur ces deux points,
nous ne pouvons pas conclure, et nous en donnons les raisons. »
\end{boxscenario}

\subsection{Pas de niveau de maturité pour un domaine opaque}

Les directions aiment résumer un domaine par un niveau de 1 à 5. Le corpus
propose une conversion du taux de conformité en niveau de maturité~:
moins de 0,35, niveau~1~; de 0,35 à 0,55, niveau~2~; de 0,55 à 0,75,
niveau~3~; de 0,75 à 0,90, niveau~4~; au-delà, niveau~5. C'est une
\emph{convention pédagogique déclarée}, non une norme, et elle a une règle~:
si $U \geq 0{,}20$, le niveau est \emph{indéterminé}. Le réseau d'agences
n'a donc pas de niveau en 2028~; lui en attribuer un, sur la foi de
$\mu = 0{,}375$, ferait croire au comité qu'il connaît un état qu'en partie
il ignore.

% ============================================================================
\section{Rédiger l'Observation}
\label{sec:c07-observation}
% ============================================================================

\subsection{Sept éléments}

Chaque non-conformité significative donne lieu à une observation, de
structure invariante.

\begin{tableaucours}{cL{3.0cm}L{8.4cm}}{%
    \tableauHeader{N°} &
    \tableauHeader{Élément} &
    \tableauHeader{Contenu}
}
1 & Norme & Le critère, et le texte ou le référentiel d'où il vient \\
2 & Constatation & Le fait établi, chiffré, avec sa preuve et sa cote \\
3 & Risques & Ce qui peut arriver~: le risque de la fiche d'assertion et sa
conséquence \\
4 & Contradictoire & La réponse de l'audité, datée \\
5 & Analyse & La cause de l'écart et l'interprétation, avec un niveau de
confiance \\
6 & Conclusion & La position finale de l'auditeur \\
7 & Recommandations & Ce qu'il faut faire, en traitant la cause \\
\end{tableaucours}
\vspace{-0.8em}
\begin{center}{\small\itshape Tableau~7.6. Les sept éléments de
l'observation}\end{center}

\begin{boxscenario}{Opération SANAGA. Observation S4.2 · Mode dégradé}
\textbf{1. Norme~:} plan de traitement de R-AG-01, action A5 («~procédure
de guichet en mode dégradé, exercée~»)~; ISO/IEC~27002:2022, mesure~5.30
(préparation des TIC pour la continuité d'activité)~\cite{iso27002_2022}.

\medskip
\textbf{2. Constatation~:} exercice inopiné du 14~mars 2028. Garoua sert
un premier client en 40~minutes. À Kribi, la procédure est introuvable, le
chef d'agence remplaçant ignore les codes du mode dégradé, et aucun client
n'est servi pendant 2~h~15 (compte rendu signé, chronométrage). Cote~:
\NON{}.

\medskip
\textbf{3. Risques~:} une agence touchée par un incident reste fermée~;
perte d'exploitation et atteinte à la réputation.

\medskip
\textbf{4. Contradictoire~:} la direction du réseau d'agences indique, le
20~mars, que la procédure a été diffusée par courriel en mai~2027 et que le
chef d'agence titulaire de Kribi était en congé.

\medskip
\textbf{5. Analyse} (confiance élevée)~: la procédure a été diffusée par
la messagerie, c'est-à-dire par l'outil qui manquera le jour d'un incident.
Elle n'a jamais été intégrée au dossier papier de l'agence, et sa
connaissance repose sur une seule personne.

\medskip
\textbf{6. Conclusion~:} le contrôle S4.2 n'est pas efficace pour
l'échantillon~; son efficacité dans les 36 autres agences n'est pas
établie.

\medskip
\textbf{7. Recommandations~:} placer la procédure et les codes, sous
enveloppe scellée, dans le coffre de chaque agence~; former au moins deux
personnes par agence~; exercer chaque agence une fois par an, avec
chronométrage~; vérification par l'audit interne lors du suivi.
\end{boxscenario}

\subsection{Le nouveau risque}

La limitation de l'étendue sur les GAB n'appelle pas d'observation au sens
strict~: on ne peut pas décrire un écart que l'on n'a pas vu. Elle appelle
une appréciation de risque, que le rapport demande au gestionnaire des
risques.

\begin{boxobjet}{R-AG-02 · Compromission d'un GAB}
\textbf{Valeur métier~:} le retrait d'espèces dans les agences.

\textbf{Biens supports~:} automates de retrait de 14 agences, sous un
système d'exploitation qui ne reçoit plus de correctifs~; contrat de
maintenance sans droit d'audit.

\textbf{Événement redouté~:} retraits frauduleux, ou indisponibilité des
automates.

\textbf{Vraisemblance~:} \NEK{} (configuration non vérifiable).

\textbf{Impact~:} à apprécier (montant des encaisses, nombre de clients
concernés).

\textbf{Action immédiate~:} négociation d'un avenant de droit d'audit avec
le prestataire~; à défaut, présentation au comité des risques, comme pour
DIGICORE (Chapitre~\ref{chap:iso27005}).
\end{boxobjet}

% ============================================================================
\section{Ce que la Conclusion demande à chaque Métier}
\label{sec:c07-metiers}
% ============================================================================

\begin{boxmetier}{Auditeur en cybersécurité}
\textbf{Conclure par sous-domaine}
\begin{enumerate}
    \item Décomposer le domaine avant le terrain, et faire revoir la
          décomposition.
    \item Rédiger une fiche d'assertion atomique pour chaque contrôle.
    \item Calculer $\mu$, $U$ et $H$ par sous-domaine, et appliquer la règle
          de décision écrite dans le plan d'audit.
    \item Conclure sur l'opacité quand $U \geq 0{,}20$, au lieu de la taire.
    \item Développer chaque non-conformité significative en observation à
          sept éléments.
\end{enumerate}
\end{boxmetier}

\begin{boxmetier}{RSSI}
\textbf{Lire le rapport par sous-domaine}

La RSSI ne retient pas « 37~\% ». Elle retient que ses barrières
techniques tiennent, qu'elle a trois défauts à corriger, et que deux
sujets lui échappent~: le personnel et les GAB. Son prochain rapport au
comité suivra la même décomposition.
\end{boxmetier}

\begin{boxmetier}{Responsable GRC}
\textbf{Une même décomposition partout}

Si le registre des risques, le plan de traitement et le programme d'audit
emploient la même décomposition (ici, les étapes du scénario plus le
facteur humain et les tiers), les cotes de l'audit alimentent directement le
registre. C'est le dispositif unique du Chapitre~\ref{chap:grc-pratique}.
\end{boxmetier}

\separateur

\begin{boxresume}
\textbf{Ce qu'il faut retenir de ce chapitre~:}
\begin{itemize}
    \item \textbf{Le gabarit canonique} a huit niveaux~: domaine,
          sous-domaine, processus, activité, contrôle, test, preuve, trace.
    \item \textbf{Une bonne décomposition} est exhaustive, non redondante et
          respecte la borne de sept éléments environ.
    \item \textbf{La fiche d'assertion} donne objectif, critère, question,
          risque, conséquence~; elle est atomique.
    \item \textbf{La cote d'un contrôle}~: un \NON{} suffit~; à défaut un
          \NEK{}~; \OUI{} seulement si tous les critères le sont.
    \item \textbf{La règle de décision}~: $U$ d'abord, $\mu$ ensuite, $H$
          pour contrôler, les cotes de l'auditeur prévalent.
    \item \textbf{Une entropie nulle} n'est pas une certitude si $U = 1$.
    \item \textbf{L'observation} a sept éléments, dont le contradictoire et
          une conclusion explicite.
\end{itemize}
\end{boxresume}

\begin{boxinfo}
\textbf{Points de contrôle.} Vous avez compris ce chapitre si vous savez~:
\begin{itemize}
    \item placer un élément au bon niveau du gabarit~;
    \item scinder une assertion qui n'est pas atomique~;
    \item calculer les trois mesures d'un sous-domaine et le situer dans le
          quadrant~;
    \item rédiger une conclusion qui dit ce qui va mal et ce que l'on ne
          sait pas.
\end{itemize}
\end{boxinfo}

% ============================================================================
\section*{Exercices}
% ============================================================================

\begin{boxexo}[title={\ding{45}\quad Exercice~7.1. À quel niveau~? \niveaubase}]
Placez chaque élément au niveau du gabarit canonique qui lui correspond.
\begin{enumerate}[(a)]
    \item L'extraction des groupes d'administration locale, avec son
          empreinte.
    \item La propagation.
    \item Le retrait des droits d'administration aux chefs d'agence.
    \item L'entrée du journal qui enregistre la modification d'un groupe
          local.
    \item Comparer les membres des groupes locaux à la liste des
          administrateurs autorisés.
    \item La gestion des droits sur les postes de travail.
    \item Aucun poste d'agence ne compte d'administrateur local autre que
          le compte central de la DSI.
\end{enumerate}
\end{boxexo}

\begin{boxexo}[title={\ding{45}\quad Exercice~7.2. Rendre atomique \niveaumoyen}]
La première version du programme portait la question suivante~: « Les
chefs d'agence ont-ils perdu leurs droits d'administration et les postes
sont-ils gérés par l'outil central ou par un prestataire agréé~? »
\begin{enumerate}[(a)]
    \item Expliquez pourquoi cette question n'admet pas de réponse honnête
          dans une agence où les droits ont été retirés mais où les postes
          ne sont gérés par personne.
    \item Scindez-la en questions atomiques.
    \item Rédigez la fiche d'assertion complète du contrôle «~fin de
          l'administration locale~» (S2.2), avec deux critères.
\end{enumerate}
\end{boxexo}

\begin{boxexo}[title={\ding{45}\quad Exercice~7.3. Le domaine D de LiZbiZ \niveaumoyen}]
AUDIX CONSEIL a décomposé le domaine~D (contrôle d'accès à l'application
intégrée de LiZbiZ) en cinq sous-domaines~:
\begin{itemize}
    \item D.1, politique et organisation~: 4 / 0 / 0~;
    \item D.2, cycle de vie des comptes~: 1 / 1 / 2~;
    \item D.3, authentification~: 4 / 0 / 0~;
    \item D.4, comptes à privilèges~: 2 / 0 / 2~;
    \item D.5, journalisation et revue~: 2 / 1 / 1.
\end{itemize}
(nombre de \OUI{}, \NON{} et \NEK{}).
\begin{enumerate}[(a)]
    \item Calculez $\mu$, $U$ et $H$ pour chaque sous-domaine et pour le
          domaine.
    \item Situez chacun dans le quadrant et appliquez la règle de décision.
    \item Quel niveau de maturité peut-on attribuer au domaine~?
    \item Où portez-vous les tests du second passage, et pourquoi~?
\end{enumerate}
\end{boxexo}

\begin{boxexo}[title={\ding{45}\quad Exercice~7.4. Observation et conclusion \niveauexpert}]
\begin{enumerate}[(a)]
    \item Rédigez l'observation à sept éléments du test T6 (journalisation~:
          31 agences sur 38, corrigée le 22~mars 2028, Chapitre~\ref{chap:audit}).
    \item Rédigez le paragraphe de conclusion du rapport consacré à S5
          (facteur humain), en concluant sur l'opacité.
    \item Un administrateur demande que le rapport «~retire les deux
          \NEK{}, qui affaiblissent le message~». Répondez-lui en citant
          les principes de l'audit (Chapitre~\ref{chap:audit}).
\end{enumerate}
\end{boxexo}

\bigskip

\begin{boxinfo}
\textbf{Pour aller plus loin.}
\begin{itemize}
    \item L'article \emph{Vers l'auditologie}~\cite{minka2026auditologie},
          pour le gabarit canonique et l'observabilité partielle.
    \item L'article sur l'incertitude épistémique~\cite{minka2026nek}, pour
          les trois mesures et le quadrant.
    \item Miller~\cite{miller1956}, sur la borne de la mémoire de travail.
\end{itemize}
\end{boxinfo}

\begin{boxinfo}
\textbf{Transition.} La Partie~II est achevée~: on sait apprécier,
traiter, chiffrer, auditer et conclure. Il reste à connaître les textes
auxquels se mesurent risques et audits~: la famille des normes ISO/IEC
27000, les autres référentiels et le droit. C'est l'objet de la
Partie~III, qui s'ouvre avec le Chapitre~\ref{chap:iso27000}.
\end{boxinfo}
